\documentclass{article}%
\usepackage[T1]{fontenc}%
\usepackage[utf8]{inputenc}%
\usepackage{lmodern}%
\usepackage{textcomp}%
\usepackage{lastpage}%
\usepackage{geometry}%
\geometry{margin=2.2cm}%
%
\usepackage{amsmath}%
\usepackage{amssymb}%
\usepackage{booktabs}%
\usepackage{longtable}%
\usepackage{array}%
\usepackage{hyperref}%
\hypersetup{hidelinks}%
\title{IME 2026 {-} Matemática (claude\_opus\_5\_5) x Poliedro\\[4pt]\large Comparação com o gabarito Poliedro Resolve\thanks{Gerado com gpt-resolve: \url{https://github.com/lgabs/gpt-resolve}}}%
\date{10/11/2026}%
%
\begin{document}%
\normalsize%
\maketitle%
\section*{Resumo}\textbf{Nota média: 10,0/10} (10 de 10 itens com a resolução correta; cada item vale o mesmo dentro da questão).\par\medskip%
\begin{longtable}{@{}c c p{0.78\textwidth}@{}}\toprule Questão & Nota & Resumo \\ \midrule
1 & 10,0 & A resolução está correta e coincide com o gabarito Poliedro: $\alpha$ = 3 e $\beta$ = 2 (z = 3 + 2i). Os caminhos são diferentes. A resolução usa o módulo ($\alpha$$^2$+$\beta$$^2$ = 13) e depois a equação da parte real. O gabarito usa divisibilidade na parte real ($\alpha$ é múltiplo de 3 e divide 9). Os dois são válidos e completos, e não encontrei nada a levar ao professor. \\
2 & 10,0 & A resolução está correta e bate com o gabarito do Poliedro: x = -7/2 ou x = -23/2. O argumento é o mesmo nos dois: a soma dos termos x\^{}4+14x\^{}3+(147/2)x\^{}2+(343/2)x+2401/16 é (x+7/2)\^{}4, e 2x+7 = 2(x+7/2), então a equação vira y\^{}3(y+8)=0. A resolução vai um pouco além do gabarito: diz que -7/2 é raiz tripla, explica por que não há outras raízes (grau 4) e substitui as duas soluções na equação original. Conferi as contas com aritmética exata e não encontrei divergência nem ponto que precise ir a um professor. \\
3 & 10,0 & A resolução está correta e dá o mesmo resultado do gabarito Poliedro: k = (1 - a)(1 + $\surd$a). O enunciado transcrito confere com o PDF original. Refiz a minimização da inclinação da secante e conferi numericamente para vários valores de a. A resolução é mais completa que o gabarito: mostra que o vértice está no domínio e explica a exclusão de a = 1/9 (nesse caso o mínimo só seria atingido pela tangente). O gabarito tem dois erros de digitação que não mudam o resultado: escreve 'k/4' no lugar de 4/k e '$\alpha$' no lugar de a. \\
4 & 10,0 & A resolução e o gabarito do Poliedro dão o mesmo conjunto de soluções, só que em formas diferentes. São as três retas pela origem, sem a origem: b = a, b = (-2-$\surd$3)a e b = (-2+$\surd$3)a, ou seja, ângulos de 45$^\circ$, 105$^\circ$ e 165$^\circ$. A resolução fatora a equação como (A-i)(B-i) = -2 e chega a tan 3$\theta$ = -1. O gabarito chega à cúbica (a-b)(a$^2$+4ab+b$^2$) = 0. As duas equações são a mesma (cos 3$\theta$ + sen 3$\theta$ = 0), e a conferência numérica confirma o resultado. A resolução está correta e completa: trata as condições de existência, só usa passos equivalentes e verifica um caso. \\
5 & 10,0 & A resolução e o gabarito estão corretos e chegam ao mesmo resultado por caminhos diferentes. A resolução prova a desigualdade pela identidade (a+b+c)(1/a+1/b+1/c) - 9 = $\Sigma$(a-b)$^2$/(ab) $\geq$ 0 e também dá o caso de igualdade, a = b = c. O gabarito do Poliedro usa a MA $\geq$ MG e só a demonstra numa observação. Conferi as duas identidades em valores racionais exatos e ambas valem. A prova da resolução é completa e elementar e mereceria nota integral. \\
6 & 10,0 & A resolução e o gabarito Poliedro dão o mesmo valor, 6 143 850, e os dois estão corretos. O enunciado confere com o PDF original. Pelos dois caminhos, $\Sigma$(3j - 5) = 2025$\cdot$3034 e $\Sigma$a\_j - 3$\cdot$2025 = 6 149 925 - 6075, o resultado é o mesmo, e a conferência em Python confirma. A resolução é mais rigorosa que o gabarito: prova que a sequência tem de ser uma PA e que a soma não depende da solução escolhida. O gabarito apresenta a PA como \textbackslash{}"uma solução possível\textbackslash{}", o que é só uma imprecisão de redação, sem efeito na nota. \\
7 & 10,0 & Questão 7: a resposta da resolução bate com o gabarito do Poliedro, 5/16 = 31,25\%. A conta de Bayes está certa (confirmei com frações exatas) e a resolução está completa e bem justificada. O único problema está no gabarito: há um erro de digitação na expressão intermediária (0,01$\cdot$0,99 no lugar de 0,1$\cdot$0,99), mas a resposta final não muda. \\
8 & 10,0 & A resolução e o gabarito do Poliedro coincidem: a única solução real é (x, y, z) = (0, 0, 0), e as duas provas estão corretas. A unicidade vem do mesmo fato, f(t) = t + log(t + $\surd$(t$^2$+1)) ser estritamente crescente, e as duas formas de concluir x = y = z são equivalentes. A resolução é mais completa que o gabarito: justifica o domínio do log, prova que f é crescente pela derivada e verifica que (0,0,0) é solução. Não há erro nos dois. O único ponto em aberto é que o enunciado não diz a base do log, o que não muda o resultado. \\
9 & 10,0 & Conferi o enunciado da questão 9 no PDF original (página 2) e ele é igual à transcrição. A resolução dá r = a($\surd$3-1)(2-$\surd$2)/2 $\approx$ 0,214a. O gabarito da Poliedro dá r = a($\surd$3-1)/(2+$\surd$2), que é a mesma expressão antes de racionalizar o denominador. As duas estão corretas e coincidem numericamente. A resolução é completa e ainda mais rigorosa que o gabarito: justifica a simetria, descarta a raiz que cai fora do lado e confere r por dois métodos. Não encontrei erros. \\
10 & 10,0 & A resolução está certa e bate com o gabarito do Poliedro: V = ($\pi$$\surd$6/108)$\cdot$a$^3$ $\approx$ 0,0713$\cdot$a$^3$, o dobro do volume da esfera inscrita. Conferi a transcrição com o PDF original e ela está fiel. Refiz a conta (r = a$\surd$6/12, razão 1/2 entre níveis, série de razão 1/2) e o resultado se confirma. A resolução ainda é mais completa que o gabarito, porque justifica que as esferas não se sobrepõem e discute se a esfera inscrita original deve entrar na soma. \\
\bottomrule\end{longtable}%
\section*{Questão 1}%
\href{https://poliedroresolve.sistemapoliedro.com.br/vestibulares/ime/2026/ime-2-fase-dia-1-23-10-2025-2-fase-dia-1/dissertativa-1-matematica-2-fase-dia-1-ime-2026}{Resolução do gabarito}\par\smallskip%
\begin{longtable}{@{}p{0.06\textwidth} p{0.33\textwidth} p{0.33\textwidth} p{0.1\textwidth} p{0.08\textwidth}@{}}\toprule Item & Resolução & Gabarito & Veredito & Certo \\ \midrule
Questão 1 ($\alpha$ e $\beta$) & $\alpha$ = 3 e $\beta$ = 2 (z = 3 + 2i) & $\alpha$ = 3 e $\beta$ = 2 (z = 3 + 2i) & igual & ambos \\
\bottomrule\end{longtable}%
\begin{itemize}%
\item%
\textbf{Questão 1 ($\alpha$ e $\beta$))} Conferi o enunciado no PDF original e a transcrição está fiel. Respostas idênticas. A resolução usa o módulo: |z|$^3$ = $\surd$2197 = 13\^{}\{3/2\}, então $\alpha$$^2$+$\beta$$^2$ = 13. Os únicos quadrados inteiros positivos que somam 13 são 4 e 9. Pela equação (I), o caso $\alpha$$^2$=4 cai fora (dá $\alpha$ = 9/23) e o caso $\alpha$$^2$=9 dá $\alpha$=3; pela (II), $\beta$=2. O gabarito parte da (I): $\alpha$ é múltiplo de 3 e divide 9, então testa $\alpha$ $\in$ \{$\pm$3, $\pm$9\}. Só $\alpha$=3 dá $\beta$ inteiro ($\beta$=$\pm$2), e a (II) fica com $\beta$=2. Os dois métodos estão corretos. Uma busca exaustiva em python com |$\alpha$|,|$\beta$| $\leq$ 50 confirmou que (3,2) é a única solução inteira. Também está certa a observação da resolução sobre as outras raízes cúbicas: (3+2i)$\cdot$e\^{}\{2$\pi$i/3\} = -(3+2$\surd$3)/2 + (3$\surd$3-2)/2$\cdot$i, que não tem partes inteiras. A verificação (3+2i)$^3$ = -9+46i também confere.%
\end{itemize}%
\section*{Questão 2}%
\href{https://poliedroresolve.sistemapoliedro.com.br/vestibulares/ime/2026/ime-2-fase-dia-1-23-10-2025-2-fase-dia-1/dissertativa-2-matematica-2-fase-dia-1-ime-2026}{Resolução do gabarito}\par\smallskip%
\begin{longtable}{@{}p{0.06\textwidth} p{0.33\textwidth} p{0.33\textwidth} p{0.1\textwidth} p{0.08\textwidth}@{}}\toprule Item & Resolução & Gabarito & Veredito & Certo \\ \midrule
Q2 & S = \{-23/2, -7/2\} (com x = -7/2 raiz tripla) & x = -7/2 ou x = -23/2 & igual & ambos \\
\bottomrule\end{longtable}%
\begin{itemize}%
\item%
\textbf{Q2)} Conferi o enunciado no PDF original e ele bate com a transcrição. A resolução e o gabarito fazem a mesma coisa: reconhecem (x+7/2)\^{}4 = x\^{}4+14x\^{}3+(147/2)x\^{}2+(343/2)x+2401/16 e reescrevem a equação como (x+7/2)\^{}4+(2x+7)\^{}3=0. A resolução usa y = x+7/2 e chega a y\^{}3(y+8)=0. O gabarito põe (2x+7)\^{}3 em evidência e chega a (2x+7)\^{}3((2x+7)/16+1)=0. Os dois caminhos dão x=-7/2 (tripla) e x=-23/2. Conferi com frações exatas em Python: o polinômio p(x) coincide com (x+7/2)\^{}3(x+23/2) em 6 pontos, o que basta para um polinômio de grau 4, e p(-7/2)=p(-23/2)=0. Como o grau é 4 e a multiplicidade total encontrada é 4, não há outras raízes, nem complexas.%
\end{itemize}%
\section*{Questão 3}%
\href{https://poliedroresolve.sistemapoliedro.com.br/vestibulares/ime/2026/ime-2-fase-dia-1-23-10-2025-2-fase-dia-1/dissertativa-3-matematica-2-fase-dia-1-ime-2026}{Resolução do gabarito}\par\smallskip%
\begin{longtable}{@{}p{0.06\textwidth} p{0.33\textwidth} p{0.33\textwidth} p{0.1\textwidth} p{0.08\textwidth}@{}}\toprule Item & Resolução & Gabarito & Veredito & Certo \\ \midrule
k em função de a & k = (1 - a)(1 + $\surd$a) = (1 - $\surd$a)(1 + $\surd$a)$^2$; inclinação mínima 4/{[}(1 - a)(1 + $\surd$a){]}, atingida em x = ((1 - $\surd$a)/2)$^2$ & k = (1 - a)(1 + $\surd$a); inclinação mínima 4/{[}(1 - $\surd$a)(1 + $\surd$a)$^2${]} & igual & ambos \\
\bottomrule\end{longtable}%
\begin{itemize}%
\item%
\textbf{k em função de a)} Refiz a conta. f(x) = $\surd$x/(1 - $\surd$x). Com s = $\surd$a e t = $\surd$x, a inclinação da secante é m = 1/{[}(1 - s)(1 - t)(t + s){]}. O produto (1 - t)(t + s) tem máximo em t$_{0}$ = (1 - s)/2, que fica em {[}0,1) e vale ((1 + s)/2)$^2$. Assim, m\_min = 4/{[}(1 - s)(1 + s)$^2${]} = 4/{[}(1 - a)(1 + $\surd$a){]}. A reta ky - 4x = 2026 tem inclinação 4/k, logo k = (1 - a)(1 + $\surd$a). Conferi numericamente, com varredura em x, para a = 0; 0,04; 0,25; 0,5; 0,81; 0,95: o k obtido e o ponto de mínimo x = ((1 - $\surd$a)/2)$^2$ coincidem com a fórmula. As duas formas de m\_min, a da resolução e a do gabarito, são idênticas, porque (1 - a) = (1 - $\surd$a)(1 + $\surd$a). A resolução também explica por que a $\neq$ 1/9: nesse caso t$_{0}$ = s, e o mínimo seria atingido só pela tangente.%
\end{itemize}%
\paragraph{Pontos de dúvida}%
\begin{itemize}%
\item%
\textit{(key)} O texto do gabarito diz que o coeficiente angular de ky - 4x = 2026 é k/4, mas o certo é 4/k. A própria conta do gabarito usa 4/k e chega ao k correto, então é só erro de digitação, sem efeito no resultado.%
\item%
\textit{(key)} O gabarito escreve o ponto como ($\alpha$, $\surd$a/(1 - $\surd$a)), com $\alpha$ em vez de a, o que parece erro de digitação ou de OCR. Ele também não justifica que o vértice t$_{0}$ = (1 - $\surd$a)/2 está em {[}0,1) nem que é diferente de $\surd$a, que é o papel da condição a $\neq$ 1/9. A resolução justifica os dois pontos.%
\end{itemize}%
\section*{Questão 4}%
\href{https://poliedroresolve.sistemapoliedro.com.br/vestibulares/ime/2026/ime-2-fase-dia-1-23-10-2025-2-fase-dia-1/dissertativa-4-matematica-2-fase-dia-1-ime-2026}{Resolução do gabarito}\par\smallskip%
\begin{longtable}{@{}p{0.06\textwidth} p{0.33\textwidth} p{0.33\textwidth} p{0.1\textwidth} p{0.08\textwidth}@{}}\toprule Item & Resolução & Gabarito & Veredito & Certo \\ \midrule
Questão 4 (valores de z) & z = r(cos $\theta$ + i sen $\theta$), com r > 0 e $\theta$ = $\pi$/4 + k$\pi$/3 (k $\in$ $\mathbb{Z}$). São três retas pela origem, sem a origem, com ângulos de 45$^\circ$, 105$^\circ$ e 165$^\circ$. Na forma algébrica: z = a + bi, a $\neq$ 0, com b = a, b = -(2+$\surd$3)a ou b = ($\surd$3-2)a. & z = a + ai ou z = a + a(-2 $\pm$ $\surd$3)i, com a $\in$ $\mathbb{R}$* & equivalente & ambos \\
\bottomrule\end{longtable}%
\begin{itemize}%
\item%
\textbf{Questão 4 (valores de z))} Os dois conjuntos são iguais. tan 45$^\circ$ = 1, tan 105$^\circ$ = -(2+$\surd$3) = -2-$\surd$3 e tan 165$^\circ$ = $\surd$3-2 = -2+$\surd$3, que são exatamente os coeficientes b/a do gabarito. Os dois caminhos levam à mesma equação. Com a = cos $\theta$ e b = sen $\theta$, a equação cúbica do gabarito a$^3$+3a$^2$b-3ab$^2$-b$^3$ = 0 vira cos 3$\theta$ + sen 3$\theta$ = 0, a mesma da resolução. Também conferi numericamente, em Python. A expressão do lado esquerdo vale 1 nas três retas, para vários valores de a, positivos e negativos. Ela não depende do módulo r. Uma varredura de $\theta$ em {[}0, 2$\pi$) acha zeros só em 45$^\circ$, 105$^\circ$, 165$^\circ$, 225$^\circ$, 285$^\circ$ e 345$^\circ$, o que confirma que não há outras soluções. As condições de existência (Re z $\neq$ 0 e Im z $\neq$ 0) valem nas duas respostas: a $\neq$ 0 e as inclinações são não nulas. O enunciado transcrito confere com o PDF original.%
\end{itemize}%
\paragraph{Pontos de dúvida}%
\begin{itemize}%
\item%
\textit{(key)} Na resolução do gabarito, algumas equações aparecem coladas na mesma linha ("... = 1 a/b + ..." e a fatoração de a$^3$+3a$^2$b-3ab$^2$-b$^3$). Isso parece problema de diagramação ou extração, não erro de conta. O gabarito também não menciona as condições de existência (Re z $\neq$ 0, Im z $\neq$ 0), mas elas ficam garantidas por a $\in$ $\mathbb{R}$*.%
\item%
\textit{(solution)} A resolução apresenta duas formas equivalentes da resposta: a polar, com r > 0 e $\theta$ = $\pi$/4 + k$\pi$/3, e a algébrica, com a $\neq$ 0 e os três coeficientes. Chama de "inclinação" o ângulo da reta, e não o coeficiente angular, mas dá os dois. É só redação: as duas formas descrevem o mesmo conjunto.%
\end{itemize}%
\section*{Questão 5}%
\href{https://poliedroresolve.sistemapoliedro.com.br/vestibulares/ime/2026/ime-2-fase-dia-1-23-10-2025-2-fase-dia-1/dissertativa-5-matematica-2-fase-dia-1-ime-2026}{Resolução do gabarito}\par\smallskip%
\begin{longtable}{@{}p{0.06\textwidth} p{0.33\textwidth} p{0.33\textwidth} p{0.1\textwidth} p{0.08\textwidth}@{}}\toprule Item & Resolução & Gabarito & Veredito & Certo \\ \midrule
Questão 5 (demonstração) & Prova direta: (a+b+c)(1/a+1/b+1/c) - 9 = (a-b)$^2$/(ab) + (a-c)$^2$/(ac) + (b-c)$^2$/(bc) $\geq$ 0, logo o produto é $\geq$ 9. A igualdade vale se e somente se a = b = c. Como prova alternativa, cita MA $\geq$ MG. & Por MA $\geq$ MG: a+b+c $\geq$ 3?(abc) e 1/a+1/b+1/c $\geq$ 3/?(abc). Como os dois lados são positivos, multiplicam{-}se as desigualdades e obtém{-}se $\geq$ 9 (c.q.d.). Em observação, o gabarito sugere demonstrar a MA $\geq$ MG para três termos via p$^3$+q$^3$+r$^3$-3pqr = (p+q+r)$\cdot$$\frac12${[}(p-q)$^2$+(p-r)$^2$+(q-r)$^2${]} $\geq$ 0. & equivalente & ambos \\
\bottomrule\end{longtable}%
\begin{itemize}%
\item%
\textbf{Questão 5 (demonstração))} As duas provas chegam à mesma conclusão por caminhos diferentes e as duas estão corretas. Conferi o enunciado no PDF original e ele está igual à transcrição. Testei as duas identidades com aritmética racional exata em 2000 trios aleatórios e ambas valem: (a+b+c)($\Sigma$1/a) - 9 = $\Sigma$(a-b)$^2$/(ab) e p$^3$+q$^3$+r$^3$-3pqr = (p+q+r)$\cdot$$\frac12$$\Sigma$(p-q)$^2$. No exemplo da resolução, (1,2,3) dá 11, como ela afirma. A prova principal da resolução é elementar e completa: só usa a distributiva e o fato de que um quadrado é $\geq$ 0, sem citar nenhum teorema. Ela também identifica corretamente o caso de igualdade, a = b = c. A prova principal do gabarito apoia{-}se na MA $\geq$ MG e só a demonstra numa observação opcional. A demonstração dessa observação também está correta.%
\end{itemize}%
\paragraph{Pontos de dúvida}%
\begin{itemize}%
\item%
\textit{(key)} A prova principal do gabarito usa a MA $\geq$ MG para três termos sem demonstrá{-}la e só sugere a demonstração numa observação. Se a banca exigir a prova da MA $\geq$ MG, a versão curta do gabarito pode perder pontos. A resolução conferida não tem esse problema, porque a prova principal dela é autossuficiente.%
\item%
\textit{(solution)} A 'Observação (outra prova)' da resolução também usa a MA $\geq$ MG sem demonstrá{-}la. Isso não afeta a nota, porque é só um comentário extra: a prova principal, pela identidade com quadrados, está completa.%
\end{itemize}%
\section*{Questão 6}%
\href{https://poliedroresolve.sistemapoliedro.com.br/vestibulares/ime/2026/ime-2-fase-dia-1-23-10-2025-2-fase-dia-1/dissertativa-6-matematica-2-fase-dia-1-ime-2026}{Resolução do gabarito}\par\smallskip%
\begin{longtable}{@{}p{0.06\textwidth} p{0.33\textwidth} p{0.33\textwidth} p{0.1\textwidth} p{0.08\textwidth}@{}}\toprule Item & Resolução & Gabarito & Veredito & Certo \\ \midrule
6 & 6143850 (= 2025$\cdot$3034); a\_n = 3n - 2, y\_j = 3, x\_j + y\_j + z\_j = a\_j - 3 = 3j - 5 & 6 143 850 (PA de razão 3, y\_j = 3, x\_j + y\_j + z\_j = a\_j - 3; soma = (1+6073)$\cdot$2025/2 - 2025$\cdot$3) & igual & ambos \\
\bottomrule\end{longtable}%
\begin{itemize}%
\item%
\textbf{6)} Conferi o enunciado com o PDF original (pág. 2) e ele bate com a transcrição. Os dois chegam ao mesmo valor. Subtraindo as equações duas a duas, obtém{-}se y = a\_\{j+1\} - a\_j = a\_\{j+2\} - a\_\{j+1\}. Como isso vale para j = 1,\ldots{},2025, todas as diferenças consecutivas de a\_1 até a\_2027 são iguais. Logo a sequência é uma PA, com r = (10 - 1)/3 = 3 e a\_j = 3j - 2. Em cada sistema, x + z = a\_j - 2y = a\_j - 6, então x\_j + y\_j + z\_j = a\_j - 3 = 3j - 5, qualquer que seja a solução escolhida. A soma é $\Sigma$(3j - 5) = 3$\cdot$2025$\cdot$2026/2 - 5$\cdot$2025 = 2025$\cdot$3034 = 6 143 850. Pelo caminho do gabarito: $\Sigma$a\_j = (1 + 6073)$\cdot$2025/2 = 6 149 925, e 6 149 925 - 6075 = 6 143 850. Conferi em Python, recalculando y e s em cada um dos 2025 sistemas. A resolução é até mais rigorosa que o gabarito: prova que a condição é necessária e suficiente e explica por que a soma não depende da solução escolhida, num sistema possível e indeterminado.%
\end{itemize}%
\paragraph{Pontos de dúvida}%
\begin{itemize}%
\item%
\textit{(key)} O gabarito diz que "uma solução possível é a sequência (a\_n) sendo uma PA", como se fosse só um dos casos. Na verdade, ser PA é obrigatório: as subtrações mostram que, se o sistema é possível para todo j, a sequência precisa ser PA. A redação é imprecisa, mas não muda o resultado.%
\item%
\textit{(key)} O gabarito não comenta que o sistema é possível e indeterminado (x e z só aparecem como x + z). Também não justifica que x\_j + y\_j + z\_j não depende da solução escolhida, o que é o que torna a soma bem definida. Ele usa x\_j + z\_j = a\_j - 6 sem destacar esse ponto. A resolução trata isso explicitamente.%
\end{itemize}%
\section*{Questão 7}%
\href{https://poliedroresolve.sistemapoliedro.com.br/vestibulares/ime/2026/ime-2-fase-dia-1-23-10-2025-2-fase-dia-1/dissertativa-7-matematica-2-fase-dia-1-ime-2026}{Resolução do gabarito}\par\smallskip%
\begin{longtable}{@{}p{0.06\textwidth} p{0.33\textwidth} p{0.33\textwidth} p{0.1\textwidth} p{0.08\textwidth}@{}}\toprule Item & Resolução & Gabarito & Veredito & Certo \\ \midrule
Questão 7 (única) & P(I | resposta 'não infectado') = 0,045/0,144 = 5/16 = 31,25\% & 5/16 & igual & ambos \\
\bottomrule\end{longtable}%
\begin{itemize}%
\item%
\textbf{Questão 7 (única))} O enunciado transcrito confere com o original (pág. 2). Refiz a conta com frações exatas: P(I)$\cdot$P(D?|I) = 0,9$\cdot$0,05 = 9/200 e P(?)$\cdot$P(D?|?) = 0,1$\cdot$0,99 = 99/1000. Logo P(I|D?) = (9/200)/(9/200 + 99/1000) = 5/16. As respostas finais são iguais. A resolução está completa: define os eventos, justifica os complementos, aplica probabilidade total e Bayes e confere que 5/16 + 11/16 = 1. A expressão escrita na resolução do gabarito tem um erro de digitação no denominador, mas o valor final está certo.%
\end{itemize}%
\paragraph{Pontos de dúvida}%
\begin{itemize}%
\item%
\textit{(key)} Na resolução do Poliedro, o denominador está escrito como 0,9$\cdot$0,05 + 0,01$\cdot$0,99, mas o certo é 0,1$\cdot$0,99, com P(não infectado) = 0,1. Com o valor 0,01 impresso, a conta daria 50/61 $\approx$ 0,82, e não 5/16. O resultado final 5/16 do gabarito está correto, então é só erro de digitação. Na transcrição também se perderam as barras de complemento: aparece 'P(B/A)' três vezes, quando deveria ser P(B|A)=0,95, P(B|?)=0,01 e P(B?|A)=0,05.%
\item%
\textit{(statement)} Lemos 'indica a presença de vírus quando o computador não está infectado com probabilidade de 1\%' como P(indica vírus | não infectado) = 0,01, a leitura usual de falso positivo. Tanto a resolução quanto o gabarito usam essa leitura, e não há outra plausível que mude o resultado.%
\end{itemize}%
\section*{Questão 8}%
\href{https://poliedroresolve.sistemapoliedro.com.br/vestibulares/ime/2026/ime-2-fase-dia-1-23-10-2025-2-fase-dia-1/dissertativa-8-matematica-2-fase-dia-1-ime-2026}{Resolução do gabarito}\par\smallskip%
\begin{longtable}{@{}p{0.06\textwidth} p{0.33\textwidth} p{0.33\textwidth} p{0.1\textwidth} p{0.08\textwidth}@{}}\toprule Item & Resolução & Gabarito & Veredito & Certo \\ \midrule
Q8 (unicidade e cálculo da solução) & Solução única (x, y, z) = (0, 0, 0). A unicidade vem de f(t) = t + log(t + $\surd$(t$^2$+1)) ser estritamente crescente, o que força x = y = z. Daí log(x + $\surd$(x$^2$+1)) = 0, logo x = 0. A existência é verificada por substituição. & Solução única (x, y, z) = (0, 0, 0). Como f é estritamente crescente, x = f(f(f(x))) implica f(x) = x, logo x = y = z. Então $\surd$(x$^2$+1) = 1 - x, que dá x = 0. & igual & ambos \\
\bottomrule\end{longtable}%
\begin{itemize}%
\item%
\textbf{Q8 (unicidade e cálculo da solução))} As duas chegam ao mesmo resultado e usam o mesmo argumento central: f é estritamente crescente. A resolução trata os casos x<y e x>y, que dão as cadeias x<y<z<x e x>y>z>x, ambas absurdas. O Poliedro compara f(x) com x dentro de f?f?f. Os dois caminhos são válidos e equivalentes. A resolução é mais completa em três pontos: mostra que o log está definido para todo real (t + $\surd$(t$^2$+1) > 0), prova que f é crescente pela derivada de h(t) = t + $\surd$(t$^2$+1) e verifica explicitamente que (0,0,0) satisfaz o sistema. O gabarito só afirma que f é crescente e não faz a verificação, mas isso é economia de redação, não erro. Conferi numericamente nas bases 10 e e: f(0) = 0, f é crescente e f(t) - t tem o mesmo sinal de t, então o único ponto fixo é 0.%
\end{itemize}%
\paragraph{Pontos de dúvida}%
\begin{itemize}%
\item%
\textit{(statement)} O enunciado não diz a base do log (decimal ou natural). Isso não muda o resultado: com qualquer base b > 1, log(t + $\surd$(t$^2$+1)) = arcsenh(t)/ln b é crescente e f(t) = t só vale em t = 0. A resolução diz isso explicitamente; o gabarito não comenta.%
\item%
\textit{(key)} O gabarito do Poliedro afirma que f é estritamente crescente sem justificar. Também não verifica que (0,0,0) satisfaz o sistema nem explica por que t + $\surd$(t$^2$+1) > 0, que é o que garante o domínio do log. Como a questão pede para 'provar', uma correção rigorosa poderia querer essas justificativas. A resolução conferida tem todas elas.%
\end{itemize}%
\section*{Questão 9}%
\href{https://poliedroresolve.sistemapoliedro.com.br/vestibulares/ime/2026/ime-2-fase-dia-1-23-10-2025-2-fase-dia-1/dissertativa-9-matematica-2-fase-dia-1-ime-2026}{Resolução do gabarito}\par\smallskip%
\begin{longtable}{@{}p{0.06\textwidth} p{0.33\textwidth} p{0.33\textwidth} p{0.1\textwidth} p{0.08\textwidth}@{}}\toprule Item & Resolução & Gabarito & Veredito & Certo \\ \midrule
Q9 --- raio do círculo inscrito no triângulo CMN & r = a($\surd$3-1)(2-$\surd$2)/2 = a(2$\surd$3+$\surd$2-$\surd$6-2)/2 $\approx$ 0,214a & r = a($\surd$3-1)/(2+$\surd$2) & equivalente & ambos \\
\bottomrule\end{longtable}%
\begin{itemize}%
\item%
\textbf{Q9 --- raio do círculo inscrito no triângulo CMN)} As duas formas são iguais: multiplicando a do gabarito por (2-$\surd$2)/(2-$\surd$2) sai a($\surd$3-1)(2-$\surd$2)/2, porque (2+$\surd$2)(2-$\surd$2)=2. Ao expandir, ($\surd$3-1)(2-$\surd$2) = 2$\surd$3-$\surd$6-2+$\surd$2. Conferi em python: as três expressões valem 0,2144127\ldots{}$\cdot$a. Também conferi as etapas. x = BM = a(2-$\surd$3) = a$\cdot$tg15$^\circ$. AM$^2$ = a$^2$+x$^2$ e MN$^2$ = 2(a-x)$^2$ dão o mesmo valor, $\approx$1,0353a. Os catetos medem a($\surd$3-1) e r = (b+c-h)/2 = S/p. A resolução ainda justifica BM = DN pelo caso cateto{-}hipotenusa e descarta a raiz 2+$\surd$3 porque ela fica fora do lado. O gabarito apenas supõe o ângulo BAM = 15$^\circ$ pela figura.%
\end{itemize}%
\paragraph{Pontos de dúvida}%
\begin{itemize}%
\item%
\textit{(key)} O gabarito escreve BM = AB$\cdot$tg15$^\circ$ sem justificar a simetria BAM = DAN = 15$^\circ$. A simetria vem da congruência ABM $\cong$ ADN, pelo caso cateto{-}hipotenusa, que o gabarito toma como evidente a partir da figura. A resolução justifica esse passo. Isso não muda o resultado.%
\item%
\textit{(key)} Na transcrição do gabarito, a fração final aparece como (a$^2$($\surd$3-1)$^2$)/(2+$\surd$2)a($\surd$3-1), sem parênteses no denominador e com os fatores 1/2 já cancelados. É só uma questão de notação, e o resultado a($\surd$3-1)/(2+$\surd$2) está correto.%
\end{itemize}%
\section*{Questão 10}%
\href{https://poliedroresolve.sistemapoliedro.com.br/vestibulares/ime/2026/ime-2-fase-dia-1-23-10-2025-2-fase-dia-1/dissertativa-10-matematica-2-fase-dia-1-ime-2026}{Resolução do gabarito}\par\smallskip%
\begin{longtable}{@{}p{0.06\textwidth} p{0.33\textwidth} p{0.33\textwidth} p{0.1\textwidth} p{0.08\textwidth}@{}}\toprule Item & Resolução & Gabarito & Veredito & Certo \\ \midrule
Volume total das esferas (questão 10) & V = ($\pi$$\surd$6/108)$\cdot$a$^3$ $\approx$ 0,0713$\cdot$a$^3$ (= 2$\cdot$V0, com V0 = $\pi$$\surd$6$\cdot$a$^3$/216) & V = 2$\cdot$v0 = ($\pi$$\surd$6/108)$\cdot$a$^3$, com v0 = $\pi$$\surd$6$\cdot$a$^3$/216 & igual & ambos \\
\bottomrule\end{longtable}%
\begin{itemize}%
\item%
\textbf{Volume total das esferas (questão 10))} Refiz as contas. H = a$\surd$6/3, r = H/4 = a$\surd$6/12, V0 = (4/3)$\pi$$\cdot$r$^3$ = $\pi$$\surd$6$\cdot$a$^3$/216. O plano tangente fica a 2r = H/2 da face oposta e corta as arestas nos pontos médios, então o tetraedro de canto tem razão de semelhança 1/2. A cada nível há 4 vezes mais esferas e cada uma tem 1/8 do volume, logo o nível soma metade do anterior. A série dá V0/(1-1/2) = 2V0 = $\pi$$\surd$6$\cdot$a$^3$/108 $\approx$ 0,07125$\cdot$a$^3$. Conferi em python: a série somada numericamente dá o mesmo valor, e a razão 2V0/V\_tetraedro $\approx$ 0,605 bate com a verificação da resolução. A resolução ainda justifica que as esferas não se sobrepõem (as de níveis seguidos só se tangenciam, e os tetraedros de um mesmo nível têm interiores disjuntos). O gabarito não faz isso, mas também não precisa.%
\end{itemize}%
\paragraph{Pontos de dúvida}%
\begin{itemize}%
\item%
\textit{(statement)} O enunciado pede o volume de "todas as esferas interiores ao tetraedro ABCD". A resolução e o gabarito contam também a esfera inscrita original (termo n = 0). Quem a excluísse acharia V0 = $\pi$$\surd$6$\cdot$a$^3$/216. Essa leitura é menos natural, e a própria resolução a comenta e descarta. Pode valer confirmar com o professor se uma banca aceitaria a outra interpretação.%
\item%
\textit{(key)} O gabarito do Poliedro não justifica que as esferas são disjuntas: cada uma só tangencia a do nível seguinte, e os tetraedros de canto têm interiores disjuntos. Sem isso, não está provado que os volumes podem ser somados. A resolução conferida cobre esse ponto (item 3). Isso não muda o resultado.%
\end{itemize}%
\end{document}